\newpage
\appendix
\section*{Supplementary Information}

\setcounter{section}{0}
\renewcommand{\thesection}{S\arabic{section}}
\renewcommand{\theequation}{S\arabic{equation}}
\setcounter{equation}{0}
\renewcommand{\thefigure}{S\arabic{figure}}
\setcounter{figure}{0}
\renewcommand{\thetable}{S\arabic{table}}
\setcounter{table}{0}

\section{Cosserat Rod Geometry and Kinematic Parameterization}

We model the human spine as a Cosserat rod, defined by a centerline curve $\mathbf{r}(s,t) \in \mathbb{R}^3$ parameterized by arc length $s \in [0, L]$. This dimensionality reduction from 3D elasticity to a 1D director theory is justified by the high aspect ratio of the vertebral column ($L/d \approx 13$). To each point $s$, we attach an orthonormal material frame $\{\mathbf{d}_1(s,t), \mathbf{d}_2(s,t), \mathbf{d}_3(s,t)\}$ describing the cross-sectional orientation~\cite{goriely2017mathematics, gazzola2018forward}. The kinematics are governed by:
\begin{align}
\mathbf{r}'(s) &= \mathbf{v}(s) = \mathbf{d}_3(s) + \boldsymbol{\varepsilon}(s), \\
\mathbf{d}_i'(s) &= \mathbf{u}(s) \times \mathbf{d}_i(s),
\end{align}
where $\mathbf{v}(s)$ is the tangent vector, $\boldsymbol{\varepsilon}(s)$ represents shear and extension strains, and $\mathbf{u}(s) = \kappa_1 \mathbf{d}_1 + \kappa_2 \mathbf{d}_2 + \tau \mathbf{d}_3$ is the Darboux vector encoding flexural curvature $\kappa_{1,2}$ and torsional twist $\tau$.

\paragraph{Boundary Conditions.} The sacral base ($s=0$) is clamped and the cranial tip ($s=L$) is free:
\begin{align}
\mathbf{r}(0) = \mathbf{0}, \quad &\mathbf{d}_i(0) = \mathbf{e}_i \quad (\text{Clamped Base}) \nonumber \\
\mathbf{n}(L) = \mathbf{0}, \quad &\mathbf{m}(L) = \mathbf{0} \quad (\text{Free Tip})
\end{align}

\section{Cosserat Force and Moment Balance}

For a static rod subject to gravity $\mathbf{f}_g = \rho A \mathbf{g}$ and active moments induced by the information field, the balance equations for internal force $\mathbf{n}$ and moment $\mathbf{m}$ are:
\begin{align}
\mathbf{n}'(s) + \mathbf{f}_g &= \mathbf{0}, \nonumber \\
\mathbf{m}'(s) + \mathbf{r}'(s) \times \mathbf{n}(s) + \mathbf{m}_{\mathrm{info}}'(s) &= \mathbf{0},
\end{align}
where $\mathbf{m}_{\mathrm{info}}$ represents the active couple generated by the information field gradient.

\section{Developmental Information Field Parameterization}

The scalar information field $I(s)$ is parameterized as a bimodal Gaussian:
\begin{equation}
I(s) = A_c \exp\left[-\frac{(s/L - s_c)^2}{2\sigma_c^2}\right] + A_l \exp\left[-\frac{(s/L - s_l)^2}{2\sigma_l^2}\right] + I_0,
\end{equation}
where $A_c = 0.5$, $s_c = 0.80$, $\sigma_c = 0.08$ (cervical lordosis), $A_l = 0.7$, $s_l = 0.25$, $\sigma_l = 0.10$ (lumbar lordosis), and $I_0 = 0.3$ (baseline).

The information content is quantified using the local Shannon entropy:
\begin{equation}
\ShannonH[I] = -\int I(s) \log I(s) \, ds,
\end{equation}
which bounds the precision of the target metric specification.

\section{Biological Metric and Information Geometry}

The central hypothesis of the IEC framework is that developmental information modifies the effective metric experienced by the rod's reference manifold. We define a biological metric:
\begin{equation}
d\ell_{\mathrm{eff}}^2 = g_{\mathrm{eff}}(s)\,ds^2 = \exp\left[2\left(\beta_1 \tilde{I}(s) + \beta_2 \frac{\partial \tilde{I}}{\partial s}\right)\right] ds^2,
\end{equation}
where $\tilde{I}$ is the normalized information field and $\beta_{1,2}$ are dimensionless coupling constants.

\section{Equilibrium Deviation and Regime Classification}

We quantify the degree of biological counter-curvature using the Kullback--Leibler divergence between the IEC-coupled geometry and the passive gravitational state:
\begin{equation}
\KLDivergence = \int_0^L \left[ \frac{1}{2} \left( \kappa_{\mathrm{IEC}}(s) - \kappa_{\mathrm{passive}}(s) \right)^2 \FisherInfo(s) \right] w(s)\, ds,
\end{equation}
where $\FisherInfo(s)$ is the local Fisher Information and $w(s) = g_{\mathrm{eff}}(s)$. The normalized deviation from the gravity-only equilibrium is:
\begin{equation}
\Dgeohat = \frac{\KLDivergence}{\KLDivergence^{\mathrm{max}}(\chi_\kappa, g)}.
\end{equation}

Regime thresholds: gravity-dominated ($\Dgeohat < 0.1$), cooperative ($0.1 \leq \Dgeohat \leq 0.3$), information-dominated ($\Dgeohat > 0.3$). These were determined by analyzing the bifurcation of the lowest eigenvalue $\lambda_0$.

\section{Metabolic Power Density and Dissipation}

The metabolic power density required to sustain counter-curvature is:
\begin{equation}
\dot{\mathcal{F}} = \int_0^L \left[ \eta_p \left|\frac{\partial \kappa}{\partial t}\right|^2 + \eta_a \left(\kappa(s) - \kappa_{\mathrm{passive}}(s)\right)^2 + \Gamma_m(s) \right] ds,
\end{equation}
where $\eta_p$ is the proprioceptive feedback dissipation coefficient, $\eta_a$ is the active moment maintenance coefficient, and $\Gamma_m(s)$ is the basal tissue maintenance cost density.

\section{Protein Dynamics and Energy Competition}

During adolescent growth, mechanosensory and growth/maintenance protein classes compete for a finite metabolic energy pool:
\begin{equation}
E_{total}(t) = E_{mechano}(t) + E_{growth}(t) + E_{metabolic}
\end{equation}
\begin{equation}
\frac{dE_{mechano}}{dt} = k_{syn,m} \cdot [Resources] - k_{deg,m} \cdot [Mechano] - \dot{\mathcal{F}}_{load}
\end{equation}
where $\dot{\mathcal{F}}_{load} \propto L^3$ represents the increasing thermodynamic load, following the strain-energy route derived in the main text ($U \sim EI\kappa^2 L$ with $I \sim L^4$, $\kappa \sim L^{-1}$). The growth protein synthesis rate is driven by GH/IGF-1 and scales with $L^2$ (surface area of nutrient diffusion), creating the scaling mismatch at the molecular level.

The vulnerability is quantified via:
\begin{equation}
\Lambda(t) = \frac{dV/dt}{\tau_{adapt}}
\end{equation}
where $dV/dt$ is volumetric growth rate and $\tau_{adapt}$ is the mechanosensory adaptation time constant.

\section{The Spinal Clock and Convective Shutdown}

The intervertebral disc possesses an autonomous circadian clock driven by the BMAL1/CLOCK transcriptional loop~\cite{dudek2017clock}. Gravity provides the synchronization signal through daily loading cycles: daytime compression forces fluid out of the nucleus pulposus, while nocturnal unloading draws fluid back via osmotic pressure~\cite{urban2004nutrition}.

Under microgravity or prolonged bed rest, this pump arrests (``Convective Shutdown''), creating hypoxia that stabilizes HIF-1$\alpha$ and upregulates catabolic enzymes (MMP1, MMP3). The resulting matrix degradation reduces $B_{\mathrm{eff}}(s)$ and widens the Energy Deficit Window.

This mechanism generates three falsifiable predictions:
\begin{enumerate}
    \item Clock gene desynchronization (Per2, Bmal1 flattening in IVD) within 72 hours of unloading, preceding structural changes.
    \item HIF-1$\alpha$ stabilization in the nucleus pulposus preceding scoliotic curvature; HIF-1$\alpha$ inhibitors should delay onset.
    \item Pulsatile axial compression at the appropriate circadian phase will preserve spinal geometry more effectively than static compression.
\end{enumerate}

\subsection{Exploratory Spinal-Jetlag Simulation}

As an exploratory extension, we incorporated circadian variation into the information-coupling term and compared an entrained baseline with a high-$\chi_\kappa$ phase-shifted condition. The phase-shifted simulation produced a 2.52-fold increase in mean simulated Cobb angle compared with the entrained baseline ($24.18^\circ$ vs.\ $9.61^\circ$; Mann--Whitney U test, $p = 2.59\times 10^{-4}$). This result is retained as a supplementary hypothesis-generating analysis only. It is not used as a core claim of the main-text model and requires independent validation against disc-clock, actigraphy, and longitudinal AIS imaging data before clinical interpretation.

\section{GPU Screening and PyElastica Refinement Workflow}

We implemented an optional NVIDIA Warp workflow to accelerate broad passive-mechanics parameter screening before selecting cases for nonlinear Cosserat-rod refinement. The screening kernel evaluates the Euler--Bernoulli cantilever tip-sag approximation, the Bio-Gravitational Number $\mathcal{B}_g$, the passive sag ratio, and a dimensionless triage score across deterministic combinations of spine length ($0.25$--$0.55$\,m), radius ($0.007$--$0.014$\,m), effective Young's modulus ($10^6$--$1.5\times10^9$\,Pa), and tissue density ($900$--$1200$\,kg/m$^3$). This workflow is computational only: it is designed to identify parameter regimes requiring full nonlinear simulation, not to generate clinical predictions.

The command
\begin{verbatim}
python scripts/experiments/experiment_nvidia_warp_beam_sweep.py \
  --n-samples 100000 \
  --device auto \
  --out-dir results/nvidia_warp_beam_sweep \
  --max-csv-rows 10000 \
  --refine-pyelastica \
  --refine-count 3 \
  --refine-elements 20 \
  --refine-final-time 0.5 \
  --refine-max-sag-ratio 2.0
\end{verbatim}
was run on an NVIDIA GB10 GPU using \texttt{warp-lang} 1.13.0. The Warp kernel evaluated 100,000 candidate parameter sets in 0.0057\,s after compilation (approximately $1.74\times10^7$ samples/s). The screened distribution had median $\mathcal{B}_g = 27.41$, minimum $\mathcal{B}_g = 0.0083$, and 0.127\% of candidates below the $\mathcal{B}_g < 0.1$ low-passive-stability threshold. The 95th percentile passive sag ratio was 0.049.

The bounded refinement stage selected high-risk rows with analytic sag ratio $\leq 2.0$ and re-ran three cases through the validated PyElastica beam setup. All three PyElastica refinements completed successfully. The refined cases had $\mathcal{B}_g \approx 0.063$ and analytic sag ratios near 2.0; their large discrepancy from the small-deflection analytic formula (maximum relative error 0.991) confirms that the Warp screen correctly identifies regimes where the analytic small-deflection approximation is no longer sufficient and nonlinear Cosserat refinement is required. These GPU-screening results are therefore reported as a reproducibility and triage tool rather than as a biological validation result.

\section{Complete Theoretical Derivations}
\label{sec:supp_theory}

\textit{This section provides the complete mathematical derivations and coupling-mechanism formulations supporting the main-text conceptual summary in Section~2. It also documents, in Section~\ref{sec:supp_hopf_rejected}, the delay differential equation and Hopf-bifurcation analysis that formed the mechanism of an earlier version of this manuscript, together with the test that rejected it.}

\subsection{The Bio-Gravitational Number and the Allometric Trap}

For a vertebral column of length $L$, flexural stiffness $EI$, mass $M$, and gravitational acceleration $g$, we define the \textbf{Bio-Gravitational Number}:
\begin{equation}
\mathcal{B}_g = \frac{EI}{M g L^2},
\label{eq:bg_number_passive_supp}
\end{equation}
a dimensionless ratio of passive flexural resistance to gravitational bending moment. Low $\mathcal{B}_g$ indicates that passive stiffness alone is insufficient and shape must be actively maintained at continuous metabolic cost --- the ``Allometric Trap''. \textbf{We report as a negative result that $\mathcal{B}_g$ does not, on the present stiffness estimates, isolate the human condition.} Computed directly from the tabulated $M$, $L$ and $EI$ (Table~\ref{tab:cross_species}), $\mathcal{B}_g$ declines monotonically with body mass: only the mouse exceeds $0.1$, and the human adult ($\mathcal{B}_g = 0.0101$) is indistinguishable from the rabbit ($0.0100$) and sits nominally \emph{above} cat, dog, horse and elephant. An earlier version of this manuscript reported a posture-based dichotomy at $\mathcal{B}_g \approx 0.1$; that dichotomy was an arithmetic artifact in the tabulated column and is withdrawn. The case for obligate active maintenance in humans therefore rests on posture and axial load path --- the human spine is loaded as a vertical column carrying axial compression rather than as a horizontal beam in bending, so its critical buckling load scales as $L^{-2}$ and is acutely sensitive to growth in height~\cite{gorman2009idiopathic} --- and not on a $\mathcal{B}_g$ threshold. We retain $\mathcal{B}_g$ as a scaling descriptor only. Consistent with this weaker reading, scaling exponents vary systematically with physiological state~\cite{glazier2022variable} and vertebral scaling exhibits size-dependent breakpoints~\cite{smith2024small}; the $EI$ estimates are single-source and of heterogeneous provenance, so the cross-species ordering is provisional.

The central problem of vertebrate morphogenesis is the translation of a one-dimensional genetic code into a robust three-dimensional geometry that can withstand environmental loads. Classical structural mechanics predicts that a flexible column clamped at the base and subjected to its own weight will buckle or sag into a monotonic C-shape~\cite{goriely2017mathematics}. Yet, the healthy human spine exhibits a complex, alternating S-shape (lordosis and kyphosis).

This anomaly indicates that the S-shape is not a passive equilibrium but is actively maintained by continuous metabolic expenditure against the gravitational field. This maintenance is naturally framed as a biological instance of \emph{optimal feedback control}~\cite{todorov2002optimal}: the nervous system minimizes a cost function combining geometric error (deviation from the target shape) and metabolic effort (ATP consumption).

\subsection{The Delay-Hopf Hypothesis: Tested and Rejected}
\label{sec:supp_hopf_rejected}

\textit{An earlier version of this framework proposed that AIS onset is triggered when growth-driven lengthening of the proprioceptive feedback loop pushes the postural controller across a Hopf bifurcation. We retain the derivation here because it was central to the previous version of this manuscript, and because we are able to reject it using the framework's own model. It is presented as a falsified alternative to the recovery ratchet of Section~2, not as a supporting mechanism.}

\paragraph{The hypothesis and its formalism.}
The active maintenance of spinal alignment relies on proprioceptive feedback from paraspinal muscles and mechanosensors, and neural conduction and central processing are not instantaneous. Established models of human postural control show that the upright state is stabilized by proportional-derivative (PD) or proportional-derivative-acceleration (PDA) feedback subject to a time delay $\tau$~\cite{milton2009time, insperger2013acceleration}. Formalizing the spine as such a controller, the corrective generalized force with feedback delay $\tau$ is
\begin{equation}
F_{\mathrm{control}}(s,t) = -K_P\, e(s,\,t-\tau) \;-\; K_D\, \frac{\partial e}{\partial t}(s,\,t-\tau),
\label{eq:control_law_supp}
\end{equation}
where $K_P$ and $K_D$ are proportional and derivative gains and the tracking error is $e(s,t) = \kappa_{\mathrm{measured}}(s,t) - \kappa_{\mathrm{target}}(s)$.

Linearizing around the straight configuration, the dynamics of the spinal column become:
\begin{equation}
\rho A \frac{\partial^2 \kappa}{\partial t^2} + \eta \frac{\partial \kappa}{\partial t} + EI \frac{\partial^4 \kappa}{\partial s^4} + K_P \kappa(s, t-\tau) + K_D \frac{\partial \kappa}{\partial t}(s, t-\tau) = 0
\label{eq:dde_spine_supp}
\end{equation}

During the adolescent growth spurt the spine elongates rapidly. Because nerve conduction velocity does not scale proportionately with rapid height increases~\cite{lang1985nerve, robinson2001nerve}, the absolute neural feedback delay $\tau$ increases. The hypothesis termed this the \textbf{Derivative Gain Trap}: derivative gains ($K_D$) that provide essential damping in childhood become destabilizing once $\tau$ crosses a critical threshold $\tau_c$, at which point the system undergoes a supercritical Hopf bifurcation~\cite{landry2005dynamics}, the upright equilibrium loses stability, and limit-cycle oscillations emerge that asymmetric loading was supposed to lock into a permanent deformity.

Assuming solutions of the form $\kappa(s,t) = A\,e^{\lambda t}\sin(n\pi s/L)$, Eq.~\ref{eq:dde_spine_supp} yields the characteristic equation
\begin{equation}
\rho A\,\lambda^{2} \;+\; \eta\,\lambda \;+\; EI\left(\frac{n\pi}{L}\right)^{4} \;+\; \left(K_P + K_D\,\lambda\right)e^{-\lambda\tau} \;=\; 0 .
\label{eq:char_eq_supp}
\end{equation}
Setting $\lambda = i\omega$ gives the Hopf boundary $\tau_c(L)$, separating damped stability from sustained oscillation. The boundary has a closed form. Writing the single-mode equation in the lumped form $I\ddot\theta + b\dot\theta - mgL\,\theta + K_P\theta(t-\tau) + K_D\dot\theta(t-\tau) = 0$ (with $I = \rho A$, $b = \eta$, $mgL \to EI(n\pi/L)^4$ absorbed into the restoring term) and substituting $\lambda = i\omega$, the real and imaginary parts are two linear equations in $(\cos\omega\tau,\ \sin\omega\tau)$; eliminating them with $\cos^2 + \sin^2 = 1$ yields an \emph{exact quadratic in} $x = \omega^2$:
\begin{equation}
I^2 x^2 + \left(2 I mgL + b^2 - K_D^2\right) x + \left(mgL^2 - K_P^2\right) = 0,
\label{eq:hopf_quadratic_supp}
\end{equation}
whose positive root gives $\omega = \sqrt{x}$ and then $\tau_c = \operatorname{atan2}(\sin\omega\tau,\ \cos\omega\tau)/\omega$ from the $2\times 2$ system. No grid search, simulation, or amplitude classifier is involved (verified symbolically; \texttt{scripts/hopf\_exact.py}). Two structural facts read off the coefficients directly: (i)~whenever $K_P > mgL$ the constant term is negative, so a positive $\omega^2$ always exists --- no gain pair is delay-unconditionally stable; and (ii)~$\tau_c(K_D)$ is non-monotonic with a single interior maximum (at the model's $K_P = 120$: $K_D^* = 12.4$, $\tau_c = 79.0$\,ms), i.e.\ the derivative-gain-trap ridge exists but is bounded --- more damping buys delay margin only up to a point.

\paragraph{The decisive test, and the rejection.}
The hypothesis makes a sharp, checkable prediction: the utilisation $\tau/\tau_c$ must \emph{rise} through the growth spurt and cross unity within the AIS onset window. We evaluated this on the model's own developmental trajectory (logistic $L$: $0.6\to1.7$\,m centred at age 12; $\tau = \tau_0 + L/v$ with $\tau_0 = 0.05$\,s; gains $K_P = 20L$, $K_D = 8L$), solving Eq.~\ref{eq:char_eq_supp} for $\tau_c(L)$ at each age. \textbf{The prediction fails, and it fails in the wrong direction.}

At the conduction velocity used in the simulation ($v = 15$\,m/s), utilisation \emph{falls monotonically} from $0.806$ at age~5 to $0.655$ at age~20. It never approaches unity; its maximum over the whole interval occurs at age~5, before the spurt. The reason is that $\tau_c$ grows faster with $L$ than $\tau$ does: lengthening the spine lowers its natural frequency $\sqrt{g/L}$ and thereby buys more delay tolerance than growth consumes. Growth does not push this system toward the Hopf boundary --- it moves the system away from it.

The rejection does not depend on the choice of $v$. The value $v = 15$\,m/s in the simulation is itself roughly a quarter of the group~Ia afferent conduction velocity cited elsewhere in this work; at a physiological $v = 55$\,m/s the result is simply more emphatic, with utilisation falling $0.546 \to 0.324$ and the dimensionless delay $\tau/\sqrt{L/g}$ decreasing by $21.1\%$ across ages 5--20. Under either parameterization there is no crossing at any age.

We therefore reject the delay-Hopf mechanism as the explanation for the clustering of AIS onset at peak height velocity, and replace it with the recovery ratchet (Section~2), whose competing-rates structure reproduces that clustering without requiring a stability boundary to be crossed. We note that this reversal is consistent with the clinical picture: affected adolescents are neurologically normal and do not present with the postural-oscillation signature a Hopf bifurcation would predict. A second, independent reason concerns the plant itself: if postural control employs an internal forward model of the loop dynamics~\cite{wolpert1998internal}, raw afferent conduction delay is largely model-compensated, and the effective $\tau$ entering Eq.~\ref{eq:char_eq_supp} is a residual after prediction --- so growth-driven lengthening cannot be assumed to walk a raw delay across $\tau_c$ in the first place.

\subsection{Information-Elasticity Coupling: Curvature and Stiffness}

The IEC field enters the rod model in two places. First, the rest curvature is offset by the gradient of the positional-information field,
\begin{equation}
\boldsymbol{\kappa}_{\mathrm{rest}}(s) = \boldsymbol{\kappa}_{\mathrm{gen}} + \chi_\kappa \nabla I(s),
\label{eq:kappa_rest_iec_supp}
\end{equation}
where $\chi_\kappa$ is a geometric coupling constant and $\boldsymbol{\kappa}_{\mathrm{gen}}$ is the baseline genetic curvature. At the molecular level this coupling relies on the structural rigidity of patterning transcription factors; AlphaFold structural analysis of critical HOX proteins (e.g., HOXC8, UniProt P31273; HOXB13, UniProt Q92826) reveals conserved, high-confidence DNA-binding domains capable of enforcing the sharp spatial identity boundaries required by our field parameterization.

\paragraph{Stiffness modulation (IEC-2).}
The information field modulates the effective bending and torsional stiffness:
\begin{equation}
B_{\mathrm{eff}}(s) = E_0 I_{\mathrm{area}} (1 + \chi_E I(s)), \qquad
C_{\mathrm{eff}}(s) = G_0 J (1 + \chi_C I(s)).
\label{eq:stiffness_modulation_supp}
\end{equation}

The total energy functional governing the rod's equilibrium is:
\begin{equation}
\mathcal{E}_{\mathrm{total}} = \int_0^L \left[ \frac{1}{2} B_{\mathrm{eff}}(\kappa - \kappa_{\mathrm{rest}})^2 + \frac{1}{2} C_{\mathrm{eff}} \tau^2 + \frac{1}{2} K_{\mathrm{eff}} \varepsilon^2 + \rho A \mathbf{r} \cdot \mathbf{g} \right] ds.
\label{eq:iec_energy_supp}
\end{equation}

Here $\varepsilon$ is axial strain and the torsional term uses the geometric torsion of the rod centreline; this is unrelated to the feedback delay $\tau$ of Section~\ref{sec:supp_hopf_rejected}, which shares the symbol by convention in the two source literatures.

\subsection{The Energy Deficit Window: NAD+ and Restoration Capacity}

The biological spine is a dissipative structure requiring continuous metabolic expenditure to maintain its sensory and structural integrity. The metabolic power required to sustain this configuration scales superlinearly with length ($L^3$ demand vs $L^2$ surface-limited supply)~\cite{urban2004nutrition}, as derived in the main text.

This scaling mismatch creates a transient \textbf{Energy Deficit Window} during the adolescent growth spurt. We propose that this metabolic deficit acts on the recovery ratchet of Section~2 by suppressing the geometric-restoration rate $k_r$ --- the machinery that recovers a growth-induced deviation before it is cemented is precisely the machinery the deficit starves.

Recent evidence demonstrates that NAD+ depletion during development causes severe spinal deformities~\cite{basse2021nampt}, while oxidative stress directly drives scoliosis in animal models~\cite{pumputis2025oxidative}. Furthermore, NAD+ insufficiency activates the SARM1 pathway, which triggers local axon degeneration~\cite{gerdts2015sarm1}. We hypothesize that the adolescent Energy Deficit Window causes localized NAD+ insufficiency which mildly compromises the integrity and metabolic throughput of the mechanosensory and motor apparatus, lowering $k_r$ relative to the growth-remodeling rate $k_g(t)$ and so moving the individual toward the ratchet regime in which deviations are cemented faster than they are recovered. We emphasise that this is a hypothesis about \emph{restoration capacity}, not about feedback delay: the delay route was tested and rejected in Section~\ref{sec:supp_hopf_rejected}.

\subsection{Molecular Grounding: The AlphaFold Anisotropy Gap (Exploratory)}

Alongside the macroscopic recovery-ratchet dynamics, we explore the molecular architecture of the mechanosensory apparatus. AlphaFold v6 structural analysis of 23 key spinal proteins reveals that demand-side mechanosensory proteins (e.g., Vimentin, PIEZO2) exhibit higher \emph{structural} anisotropy $\mathcal{A}$ (mean $3.08 \pm 1.44$) than supply-side metabolic regulators (mean $1.79 \pm 0.56$).

\textbf{Symbol caution.} Two distinct quantities are called ``anisotropy'' in this literature and must not be conflated. The \emph{stiffness ratio} used in the mechanical sweeps is a directional-stiffness quantity for which a high value confers stability. The \emph{structural} anisotropy $\mathcal{A}$ defined here is a property of the proteome that enters the model as a multiplier on maintenance \emph{demand} ($P_{\mathrm{eff}} = P_{\mathrm{counter}}\,\mathcal{A}/\bar{\mathcal{A}}$), so a high value is costly, not protective: raising $\mathcal{A}$ lowers the critical spine length $L_{\mathrm{crit}}$ ($r = -0.88$ over a deterministic sweep). An earlier version of this manuscript described this relationship as ``protective scaling'', which inverted the sign; the therapeutic direction implied by the model is \emph{reduction} of $\mathcal{A}$.

We tentatively propose a \textbf{Demand-Supply Anisotropy Gap} ($p = 0.021$ two-sided, $p = 0.011$ one-sided, $n = 23$ in our dataset): the proteins required to \emph{sense} mechanical geometry are structurally extended and thermodynamically expensive to maintain, while their metabolic regulators are compact and often intrinsically disordered. We emphasize that this observation is hypothesis-generating. Current AI structural prediction tools cannot directly infer mechanical properties under load~\cite{gomes2022protein, zonta2024sequence}, and confidence scores for extracellular matrix components are frequently low. Nonetheless, this structural asymmetry provides a compelling direction for future experimental validation of the cellular mechanisms underlying the Energy Deficit Window.


\section{Detailed Methods}
\label{sec:supp_methods}

\textit{This section provides detailed computational protocols, steered molecular dynamics parameters, tissue homogenization calculations, and protein dynamics modeling supporting the main-text Methods summary in Section 3.}

\subsection{Multi-Scale Protein Mechanics Pipeline}

To ground the macroscopic parameters $\eta_p, \eta_a, \Gamma_m$ in molecular reality, we developed a multi-scale pipeline integrating AlphaFold predictions with structural energetics.
\begin{enumerate}
    \item \textbf{Structure Prediction:} We retrieved high-confidence structures for 23 key spinal proteins from the AlphaFold Protein Structure Database.
    \item \textbf{Structural Metrics:} For each protein, we computed the \emph{Anisotropy Index} (ratio of principal moments of inertia) and \emph{Radius of Gyration} ($R_g$) using the \texttt{afcc} module. These geometric metrics are reported as exploratory structural descriptors of the AlphaFold-predicted static conformation, not as direct measurements of mechanical properties.
    \item \textbf{Working hypothesis on metabolic cost:} As an exploratory hypothesis, we conjecture that the maintenance cost of a protein scales with $\text{Anisotropy} \times N_{\text{residues}}$ (extended structures plausibly require more chaperone intervention and proteostatic support). This is a working assumption, not a measured quantity, and is not load-bearing for the central control-theoretic argument; we use it only to motivate the demand-vs-supply categorization tested in Results.
    \item \textbf{Steered Molecular Dynamics (planned, not in current submission):} A subsequent work in preparation will report SMD simulations (GROMACS, CHARMM36m, TIP3P) extracting force-extension curves and single-molecule stiffness $k_{\text{molecule}}$ for the demand- and supply-side proteins, enabling direct comparison to published AFM nanoindentation data for PIEZO2~\cite{coste2010piezo}, vimentin, and lamin A/C. Quantitative validation of the anisotropy-cost mapping against single-molecule force spectroscopy is identified as the principal external validation step required before the molecular results in this paper can be elevated from exploratory to confirmatory.
\end{enumerate}

\subsection{Protein Dynamics and Energy Modeling}

We simulated the temporal competition between protein classes using a system of coupled differential equations. The model tracks the energy pool $E(t)$ available for protein synthesis/maintenance over the 5--20 year age range.
\begin{itemize}
    \item \textbf{Supply:} Modeled as $S(t) = S_0 (L(t)/L_0)^2$, reflecting the diffusion-limited transport through the vertebral endplate surface area.
    \item \textbf{Demand:} Modeled as $D(t) = D_{maint} (L/L_0) + D_{grav} (L/L_0)^4$, capturing the linear scaling of tissue volume maintenance and the quartic scaling of gravitational moment opposition.
    \item \textbf{Dynamics:} The fidelity of the mechanosensory array $F(t)$ evolves according to $dF/dt = k_{repair}(S - D)$ when $S > D$, and $dF/dt = k_{damage}(S - D)$ when $S < D$ (deficit), bounded to $[0, 1]$.
\end{itemize}


\subsection{Supplementary Table S2: Multiple-comparison correction across reported protein-statistic tests}

\begin{table}[h!]
\centering
\caption{Benjamini--Hochberg false-discovery-rate adjusted q-values (and Bonferroni-corrected p-values, for reference) across the seven independent statistical tests reported in the protein-stratification analyses (Results §3.4). Tests ordered by raw p-value. Three tests survive the standard $q < 0.05$ FDR threshold (hinge density, anisotropy gap extended trimmed, anisotropy gap original --- the last only marginally, $q = 0.049$, and only when the anisotropy-gap test is scored two-sided from the regenerated v6 table; the biped/quadruped Mann--Whitney result is reported with the explicit caveat that it does not survive correction).}
\label{tab:S2}
\small
\begin{tabular}{lcccc}
\toprule
\textbf{Test} & \textbf{N} & \textbf{$p_{\text{raw}}$} & \textbf{$p_{\text{Bonferroni}}$} & \textbf{$q_{\text{BH}}$} \\
\midrule
Mechanosensor hinge density (n=61) & 61 & 0.0025 & 0.018 & \textbf{0.018} \\
Anisotropy gap, extended cohort, trimmed & 26 & 0.0094 & 0.066 & \textbf{0.033} \\
Anisotropy gap, original demand/supply & 23 & 0.0210 (two-sided; $0.0105$ one-sided) & 0.147 & \textbf{0.049} \\
Biped vs Quadruped MWU (per-species) & 12 & 0.0440 & 0.308 & 0.077 \\
Disorder $\times$ anisotropy in mechanosensors & 31 & 0.0720 & 0.504 & 0.101 \\
AIS-strict-tagged vs other (text-match) & 61 & 0.2340 & 1.000 & 0.273 \\
AIS-strict-tagged trimmed & 59 & 0.2870 & 1.000 & 0.287 \\
\bottomrule
\end{tabular}
\end{table}

Three additional robustness analyses (correlations \#6, \#7, \#8 in the analysis pipeline at \texttt{scripts/}) provide further accountability:

\begin{itemize}
\item \textbf{Citation-bias check (correlation \#6):} no correlation between curated clinical-evidence priority score (0--100) and AlphaFold anisotropy across $n=61$ genes (Spearman $\rho = 0.016$, $p = 0.90$). The structural signal is mechanism-specific to the functional categorization, not a literature-popularity artifact.
\item \textbf{Leave-one-out robustness (correlation \#7):} for the demand-vs-supply MWU (baseline $p = 0.035$), $84.6\%$ of LOO iterations remain $p < 0.05$ (LOO median $p = 0.041$, max $p = 0.092$). Result is borderline-fragile to single supply-side gene removal --- reflects the small supply cohort ($n = 7$). For the hinge-density test (baseline $p = 0.0025$), $100\%$ of LOO iterations remain $p < 0.01$ (LOO median $p = 0.003$, max $p = 0.0042$) --- fully robust.
\item \textbf{Demand-panel expansion sensitivity (correlation \#8):} extending the demand panel from $n=19$ curated extended-filament/channel proteins to all $n=37$ mechanotransduction-tagged proteins dilutes the gap (relative-gap $93\% \rightarrow 42\%$, $p = 0.035 \rightarrow 0.33$). The published gap is therefore a property of extended-filament/channel mechanosensors specifically, not all mechanotransduction-tagged proteins. The headline numbers (72\%, $p = 0.021$ two-sided / $0.011$ one-sided, Cohen's $d = 1.13$) cited in Results §3.4 reference the curated $n = 23$ subset; the narrowed scope is reflected in the abstract and discussion.
\end{itemize}

\section{Bio-Gravitational Critical-Point Falsification (GPU Validation)}

A dedicated falsification test asked whether counter-curvature efficiency $\eta_{CC}(\mathcal{B}_g)$ exhibits a sharp phase transition at $\mathcal{B}_g^* \approx 1.0$, as originally hypothesized in early drafts. We executed 720 GPU-accelerated rod simulations (JAX, 3 scales $\times$ 8 seeds $\times$ 30 $\chi_M$ values) plus an extended sweep of 1920 simulations spanning $\mathcal{B}_g \in [0.003, 8000]$. \textbf{Result: falsified.} Sigmoid fits did not yield a universal $\mathcal{B}_g^* \approx 1$; the global median fitted critical point was $\mathcal{B}_g^* \approx 626$, with significant scale dependence (Kruskal--Wallis $p = 0.0001$). What the main text retains is the smooth, monotonic $\eta_{CC}$--$\chi_M$ scaling at human $\mathcal{B}_g$ values (Results §3.1); no $\mathcal{B}_g$ threshold is claimed, neither a sharp universal transition at $\mathcal{B}_g = 1$ nor a posture-based threshold at $0.1$ (reported in Results §3.1 as a negative result).
